\section{Билет 58. Распределение Больцмана. Распределение Максвелла-Больцмана.}

\begin{definition}
\textbf{Распределение Больцмана} обобщает барометрическую формулу на случай любого внешнего потенциального поля. Если молекула в точке с радиус-вектором $\vec{r}$ обладает потенциальной энергией $\varepsilon_p(\vec{r})$, то концентрация молекул в этой точке при температуре $T$ определяется формулой:
\begin{equation}
    n(\vec{r}) = n_0 \exp\left(-\frac{\varepsilon_p(\vec{r})}{kT}\right), \label{eq:boltzmann_dist}
\end{equation}
где $n_0$ --- концентрация в точках, где $\varepsilon_p = 0$. Распределение Больцмана показывает, что частицы стремятся занимать состояния с минимальной потенциальной энергией, но тепловое движение ``размазывает'' их по пространству.
\end{definition}

\begin{theorem}[Распределение Максвелла-Больцмана]
Для полного описания состояния идеального газа в потенциальном поле необходимо учитывать не только координаты молекул, но и их скорости. В фазовом пространстве $(\vec{r}, \vec{v})$ число молекул $dN$, находящихся в элементе объёма $dV$ около точки $\vec{r}$ и имеющих скорости в интервале $d^3v = dv_x dv_y dv_z$ около $\vec{v}$, задаётся \textbf{распределением Максвелла-Больцмана}:
\begin{equation}
    dN(\vec{r}, \vec{v}) = n_0 \left( \frac{m}{2\pi kT} \right)^{3/2} \exp\left(-\frac{\frac{mv^2}{2} + \varepsilon_p(\vec{r})}{kT}\right) dV \, d^3v. \label{eq:maxwell_boltzmann}
\end{equation}
Здесь:
\begin{itemize}
    \item $\left( \frac{m}{2\pi kT} \right)^{3/2}$ --- нормировочный множитель, обеспечивающий $\int f(\vec{v}) d^3v = 1$;
    \item $\exp\left(-\frac{mv^2}{2kT}\right)$ --- больцмановский множитель для кинетической энергии (распределение Максвелла);
    \item $\exp\left(-\frac{\varepsilon_p}{kT}\right)$ --- пространственный больцмановский множитель.
\end{itemize}
\end{theorem}

\begin{property}[Независимость координат и скоростей]
Распределение \eqref{eq:maxwell_boltzmann} факторизуется на произведение двух независимых функций:
\begin{equation}
    f(\vec{r}, \vec{v}) = f_{\text{прост}}(\vec{r}) \cdot f_{\text{скор}}(\vec{v}),
\end{equation}
где $f_{\text{скор}}(\vec{v})$ --- распределение Максвелла по скоростям, не зависящее от координат, а $f_{\text{прост}}(\vec{r})$ --- распределение Больцмана по координатам, не зависящее от скоростей. Это означает, что в равновесном состоянии скорость молекулы не зависит от того, в какой точке пространства она находится, и наоборот.
\end{property}

\begin{remark}
\textbf{Общность закона:}
\begin{enumerate}
    \item Распределение Максвелла-Больцмана справедливо для любых классических частиц, находящихся в состоянии хаотического теплового движения в консервативном поле сил (гравитационном, электростатическом, центральном и т.д.).
    \item При $\varepsilon_p = 0$ (отсутствие внешних полей) пространственный множитель равен 1, и распределение сводится к однородному распределению Максвелла.
    \item Закон является фундаментальным результатом классической статистической физики и лежит в основе расчётов атмосферного давления, распределения заряженных частиц в плазме, осмоса и многих других явлений.
\end{enumerate}
\end{remark}
